\subsection*{10. Epilogue}
The epilogue on pp.~50--51 connects the program to the CNRS application
announced in \S1. Following the teaching appraisal of \S9, Grothendieck
states that this university activity has ceased to be useful to him or
to others, and explains his request for secondment, the competing claims
on scarce research posts, and the university position that his departure
would release. The application itself has prompted the present sketch;
the table of contents was added afterward to help readers, while the
remarks about ten or a hundred years of work measure the scope he
perceives in its themes. On p.~51 he states their unity through the
geometry of surfaces, bringing together the tower and maps of \S\S2--3,
the geometric constructions of \S4, and the later foundational and
teaching discussions. His final retrospective change is a renewal of
approach to ``motifs'', whose fascination he says has persisted from his
earlier mathematical work; the footnote directs the reader to the
``divagations motiviques'' at the end of the 1972 \emph{Esquisse
Thématique}, pp.~17--18. The archive catalogue also names
\href{https://grothendieck.umontpellier.fr/archives-grothendieck/}{\emph{Motifs et théorie de Hodge} (cote 18)}.
The closing sentences describe his release from the self-imposed burden
of endless ``tâches'' and the greater lightness he associates with age.
The main report ends with the date January 1984; the subsequently added
numbered notes follow it.
(\href{https://grothendieck.umontpellier.fr/web/viewer.html?file=../119.pdf\#page=53}{Copy 1, pp.~50--51, beginning at site 52};
\href{https://grothendieck.umontpellier.fr/web/viewer.html?file=../119.pdf\#page=54}{motives, retrospective footnote, and closing date on p.~51, site 53}.)
